%!TeX program = lualatex
\renewcommand{\thelektion}{L11 \textperiodcentered{} Aufbau eines Computers II}

\begin{frame}\lessontitle{Lektion 11}{Aufbau eines Computers II}{RAM, Massenspeicher \& der Fetch-Decode-Execute-Zyklus \textperiodcentered{} Skript, Kapitel 5}\end{frame}

\begin{frame}\FDUtitle{Der Arbeitsspeicher (RAM)}
\begin{center}
Stellen Sie sich RAM als eine lange Tabelle vor: jede Zeile ist eine Speicherzelle mit eindeutiger Adresse.
\end{center}
\begin{table}[H]
    \centering
    \begin{tabular}{c|c}
        \textbf{Adresse} & \textbf{Inhalt (1 Byte)} \\\midrule\midrule
        \texttt{0x00} & \texttt{10110100} \\\midrule
        \texttt{0x01} & \texttt{00001111} \\\midrule
        \texttt{0x02} & \texttt{11001010} \\\midrule
        $\vdots$ & $\vdots$ \\
    \end{tabular}
\end{table}
\begin{center}
Die Adressen sind hexadezimal (Lektion 3) -- daher das Präfix \texttt{0x}.
\end{center}
\end{frame}

\begin{frame}\FDUtitle{RAM vs.\ Cache}
\begin{itemize}\setlength\itemsep{7pt}
  \item RAM ist \tib{flüchtig}: Inhalt geht beim Ausschalten verloren
  \item Typische Grösse: 8--64 GB
  \item Der \tib{Cache}: noch schneller, aber viel kleiner -- speichert häufig genutzte Daten zwischen
\end{itemize}
\end{frame}

\begin{frame}\FDUtitle{Massenspeicher}
\begin{columns}[T,onlytextwidth]
\begin{column}{.46\textwidth}
\large\bfseries\color{FDUteal}HDD
\vspace{2mm}
\normalsize Magnetische Scheiben speichern Bits als winzige magnetisierte Bereiche. Langsam, aber billig und langlebig.
\end{column}
\begin{column}{.46\textwidth}
\large\bfseries\color{FDUblue}SSD
\vspace{2mm}
\normalsize Daten als elektrische Ladung in Flashspeicher-Zellen. Deutlich schneller, keine beweglichen Teile.
\end{column}
\end{columns}
\vspace{5mm}
\begin{center}
RAM: 30--50 GB/s \quad SSD: 1--7 GB/s \quad HDD: 100--200 MB/s
\end{center}
\end{frame}

\begin{frame}\FDUtitle{Übung: Ladezeit}
\begin{center}
Ein Foto (24 MB) wird von einer SSD mit 2 GB/s Lesegeschwindigkeit geladen. Wie lange dauert das?
\end{center}
\only<2->{\begin{center}\Large $\dfrac{24\,\mathrm{MB}}{2000\,\mathrm{MB/s}} = 12\,\mathrm{ms}$\end{center}}
\end{frame}

{\setbeamertemplate{footline}{\darkfootline}
\begin{frame}\sectionframe[FDUgreen]{Fetch -- Decode -- Execute}{Wie die CPU Befehle abarbeitet}
\end{frame}}

\begin{frame}\FDUtitle{Der Fetch-Decode-Execute-Zyklus}
\begin{center}
\begin{tikzpicture}[
    block/.style={draw, rounded corners, minimum width=3.2cm, minimum height=1.0cm, text centered, text width=3.2cm, fill=FDUblue!15},
    arrow/.style={->, thick}
]
    \node[block] (fetch)   at (0, 0)    {\textbf{1.\ Fetch} \newline \small Befehl aus RAM holen};
    \node[block] (decode)  at (4.6, 0)  {\textbf{2.\ Decode} \newline \small Befehl entschlüsseln};
    \node[block] (execute) at (2.3,-1.9) {\textbf{3.\ Execute} \newline \small Befehl ausführen};
    \draw[arrow] (fetch)   -- (decode);
    \draw[arrow] (decode)  -- (execute);
    \draw[arrow] (execute.west) -- ++(-2.1,0) |- (fetch.south);
\end{tikzpicture}
\end{center}
\begin{center}Die CPU wiederholt diesen Ablauf Milliarden Mal pro Sekunde.\end{center}
\end{frame}

\begin{frame}\FDUtitle{Was passiert in jeder Phase?}
\begin{itemize}[<+->]\setlength\itemsep{6pt}
  \item \tib{Fetch}: Steuerwerk liest den nächsten Befehl aus dem RAM (Adresse = Programmzähler)
  \item \tib{Decode}: Befehl wird zerlegt -- \emph{was} tun, \emph{mit welchen} Operanden?
  \item \tib{Execute}: ALU (o.\,ä.) führt die Anweisung aus, Ergebnis wird gespeichert, Programmzähler rückt weiter
\end{itemize}
\end{frame}

\begin{frame}\FDUtitle{Übung: Taktfrequenz}
\begin{center}
Ein Prozessor hat $3{,}6\,\mathrm{GHz}$ Taktfrequenz und führt im Schnitt 2 Befehle pro Taktzyklus aus.
\vspace{3mm}

Wie viele Befehle führt er pro Sekunde aus? Wie lange dauert ein Programm mit $7{,}2\times 10^9$ Befehlen?
\end{center}
\only<2->{\begin{center}\Large $7{,}2\times 10^9$ Befehle/s $\quad\rightarrow\quad 1\,\mathrm{s}$\end{center}}
\end{frame}

\begin{frame}{Auftrag}{Kapitel 5: Speicher \& FDE-Zyklus}
\aufgabebox{\textbf{Aufgabe 5.2}: Dauer berechnen: Foto ($24\,\mathrm{MB}$) von SSD ($2\,\mathrm{GB/s}$) laden}\\[3mm]
\aufgabebox{\textbf{Aufgabe 5.3}: Befehle/s und Ausführungsdauer bei $3{,}6\,\mathrm{GHz}$ berechnen}\\[3mm]
\challengebox{\textbf{Aufgabe 5.4 (Challenge)}: Teilschritte des Befehls \texttt{ADD R1, R2} den Phasen Fetch, Decode, Execute zuordnen}
\end{frame}
